% BHU PYQ -- Professional Exam Template
\documentclass[11pt,a4paper]{article}

\usepackage[a4paper,top=2.5cm,bottom=2.5cm,left=2.8cm,right=2.8cm,headheight=14pt]{geometry}
\usepackage{amsmath,amssymb}
\usepackage[T1]{fontenc}
\usepackage[utf8]{inputenc}
\usepackage{lmodern}
\usepackage{microtype}
\usepackage{fancyhdr}
\usepackage{enumitem}
\usepackage{listings}
\usepackage{hyperref}
\usepackage{lastpage}
\usepackage{parskip}

\hypersetup{colorlinks=true,urlcolor=black,linkcolor=black,
  pdftitle={BVCA-305: Network Configuration and Management -- BHU PYQ}}

\pagestyle{fancy}
\fancyhf{}
\renewcommand{\headrulewidth}{0.4pt}
\renewcommand{\footrulewidth}{0.4pt}
\fancyhead[L]{\small\textit{Banaras Hindu University}}
\fancyhead[R]{\small\textit{BVCA-305: Network Configuration and Management}}
\fancyfoot[C]{\small Page \thepage\ of \pageref{LastPage}}

\lstset{
  basicstyle=\small\ttfamily,
  frame=single,
  breaklines=true,
  columns=flexible,
  keepspaces=true
}

\newlist{parts}{enumerate}{2}
\setlist[parts,1]{label=(\alph*),leftmargin=*,itemsep=4pt,topsep=3pt}
\setlist[parts,2]{label=(\roman*),leftmargin=*,itemsep=2pt,topsep=2pt}
\newcommand{\pts}[1]{\hfill\textit{[#1]}}

\begin{document}
\begin{center}
  {\large\bfseries BANARAS HINDU UNIVERSITY}\\[2pt]
  {\normalsize B.Voc. Semester III Examination, 2022-23}\\[6pt]
  \rule{0.8\linewidth}{0.5pt}\\[6pt]
  {\large\bfseries Paper: BVCA-305 --- Network Configuration and Management}\\[4pt]
  {\normalsize Time: Three hours \hfill Maximum Marks: 70}
\end{center}

\rule{\linewidth}{0.4pt}

\smallskip
\noindent\textit{Instructions: Answer all sections. Candidates are required to write their answers in their own words as far as practicable. Symbols have their usual meanings.}

\bigskip

\noindent{\large\bfseries SECTION --- A}
\rule{\linewidth}{0.4pt}\smallskip

Long Answer questions. Answer two questions, selecting one from each pair of OR questions: \hfill \textbf{2 $\times$ 15 = 30 Marks}

\medskip
\noindent\textbf{Question 1.} Describe the OSI model with a suitable diagram and compare the OSI model with the TCP/IP model.
\centerline{\textbf{OR}}
Explain the process of inter-VLAN network communication. What do you understand by a Layer 3 switch?

\medskip
\noindent\textbf{Question 2.} What is the IEEE 802.11 protocol? Explain the 802.11 infrastructure network architecture and its components.
\centerline{\textbf{OR}}
Discuss any three Internet applications with illustrative diagrams.

\bigskip\noindent{\large\bfseries SECTION --- B}
\rule{\linewidth}{0.4pt}\smallskip

\noindent\textbf{Question 3.} Attempt any six parts (Answer in about 200 words each): \hfill \textbf{6 $\times$ 5 = 30 Marks}
\begin{parts}
  \item What are the various network topologies? Discuss briefly.
  \item Describe Repeater, Hub, and Switch. Compare a Switch with a Bridge.
  \item Discuss the IEEE 802.3 Ethernet protocol and its different categories.
  \item What do you understand by a Distribution System (DS) in wireless networking?
  \item What is an Access Point (AP)? What are its various deployment methods?
  \item Explain the BGP routing protocol and discuss its two types.
  \item What is a DNS server? Explain its working with a suitable diagram.
  \item How is Wireshark helpful in network troubleshooting and packets analysis?
\end{parts}

\bigskip\noindent{\large\bfseries SECTION --- C}
\rule{\linewidth}{0.4pt}\smallskip

\noindent\textbf{Question 4.} Attempt all parts: \hfill \textbf{10 $\times$ 1 = 10 Marks}
\begin{parts}
  \item UTP stands for \_\_\_\_\_\_\_\_.
  \item In network models, cables are considered to function at Layer \_\_\_\_\_\_\_\_.
  \item In network models, routers are considered to operate at Layer \_\_\_\_\_\_\_\_.
  \item The WLC and APs use the \_\_\_\_\_\_\_\_ protocol to communicate.
  \item WPA stands for \_\_\_\_\_\_\_\_.
  \item The broadcasting address of a network like \texttt{192.168.1.0/24} is \texttt{192.168.1.}\_\_\_\_\_\_\_\_.
  \item Subnetting is the short form of \_\_\_\_\_\_\_\_.
  \item Is UTM a firewall or not?
  \item Which protocol is more reliable for data transmission --- TCP or UDP?
  \item Which command/utility is used to trace the routing path of a packet?
\end{parts}


\vfill
\rule{\linewidth}{0.4pt}
\begin{center}\small
\href{https://bhu-exam.vercel.app/}{Click here to visit our website}\\[4pt]\href{https://me-aryan.vercel.app/}{Click here to visit our contributor}\\[4pt]\href{https://quashan-me.vercel.app/}{Click here to visit our contributor}
\end{center}
\end{document}
